% d_commands.tex

% =====================================
% ---- Orthographic Commands & Others ----
% =====================================

% Commands for IPA font and IPA embedded in Japanese text
\newcommand\ip[1]{{\ipafont #1}}
\newcommand\ipe[1]{{\space\ipafont #1\space}} % % Embedded with japanese text at both side
\newcommand\ipr[1]{{\ipafont #1\space}} % Embedded with japanese text at right side
\newcommand\ipl[1]{{\space\ipafont #1}} % Embedded with japanese text at left side

% Command for embedded Ryukyuan in Japanese
%\newcommand\emb[1]{{\EmbJpFont\addfontfeatures{LetterSpace=1.0} #1}}
\newcommand\emb[1]{{\EmbJpFont #1}}

% Command for outputting roman number
\makeatletter
\newcommand*{\rom}[1]{\expandafter\@slowromancap\romannumeral #1@類}
\makeatother

% Definition of space for wakachigaki 
% NOTE: never use 全角スペース as such
%\newcommand{\ws}{\nobreak\hskip0.75em plus 0.1em minus 0.08em\relax}
\newcommand{\ws}{\space\kern0.35em}

% Kerneling of 「゜」
%\newcommand{\m}{\nobreak\hspace{-2pt}゜\hspace{-5pt plus 0.5pt minus 0.5pt}}
\newcommand{\m}{\nobreak\kern-2pt゜\kern-5pt}

\newcommand{\sms}{\kern0.05em\raisebox{-0.3ex}{\scalebox{0.7}{ス}}\kern0.05em}
\newcommand{\smz}{\kern0.05em\raisebox{-0.3ex}{\scalebox{0.7}{ズ}}\kern0.05em}

% Definition for the correct display of 「～」
\newcommand{\ntilde}{\nobreak\hspace{1pt plus 0.2pt minus 0.2pt}～\hspace{0.6pt plus 0.1pt minus 0.1pt}}

% Commands for pitch variation diacritics
\newcommand{\josho}{\pitchvar{\hspace{0pt}⸢\hspace{0pt}}}
\newcommand{\kako}{\pitchvar{\hspace{0pt}⸣\hspace{0pt}}}
\newcommand{\heiban}{\ip{\hspace{0pt}\textbf{\nolinebreak ¯}\hspace{0pt}}}

\DeclareRobustCommand{\tl}[1]{\nobreak\kern.02em\hbox{\ip{#1}}}

% Command for ①②③... 
% NOTE: used in MEANING(S) numbering
\newcommand{\numbering}[1]{%
  \ifcase#1\or ①\or ②\or ③\or ④\or ⑤\or ⑥\or ⑦\or ⑧\or ⑨\or ⑩\or ⑪\or ⑫\or ⑬\or ⑭\or ⑮\or ⑯\or ⑰\or ⑱\or ⑲\or ⑳\or ㉑\or ㉒\or ㉓\or ㉔\or ㉕\or ㉖\or ㉗\or ㉘\or ㉙\or ㉚\fi
  \hspace{0.3em}%
}

% Command for letter heading
% Principle: unbreakable vertical box with fixed height

\makeatletter
\newcommand{\LetterHead}[1]{%
  \par\noindent

  % Build an unbreakable box with *exact* spacing (tune these two):
  \setbox0=\vbox{%
    \vskip 1.0\baselineskip
    \hbox to \linewidth{\hfil \Large\bfseries —\,#1\,—\hfil}%
    \vskip 1.0\baselineskip
  }%

  % Decide remaining space (multicols column vs normal page)
  \dimen0=\ht0 \advance\dimen0 by \dp0 % needed height

  \ifdefined\@colroom
    % Inside multicols: remaining space in current column is \@colroom
    \ifdim\@colroom<\dimen0
      \columnbreak
    \fi
  \else
    % Normal single-column: remaining space on page
    \dimen2=\dimexpr\pagegoal-\pagetotal\relax
    \ifdim\dimen2<\dimen0
      \pagebreak
    \fi
  \fi

  % Place the box (unbreakable, predictable)
  \nointerlineskip\box0\par
}
\makeatother

% =====================================
% ---- DICTIONARY ENTRY STRUCTURE ----
% =====================================

\usepackage{keyval}

% New entry command
% 1st param: Entry head
% 2st param: Meanings block
% 3st param: Entry tail
\newcommand{\entry}[3]{%
	\par\begingroup
    \leavevmode
	\hangindent=0.3cm
	#1% 
	%\hspace{0.2em plus 0.1em minus 0.05em}%
	\kern0.08em%
    \ignorespaces#2#3%
    \par\endgroup
}

% Keyed Entry head
\makeatletter

% Slots for Entry head
\newcommand{\ryu@head@kana}{} % Kana orthograpy
\newcommand{\ryu@head@ipa}{} % IPA/Romaji orthograpghy
\newcommand{\ryu@head@pos}{} % Part of Speech
\newcommand{\ryu@head@cat}{} % Semantic Category 
\newcommand{\ryu@head@acc}{} % Accent Class/information
\newcommand{\ryu@head@rui}{} % Conjugation Class
\newcommand{\ryu@head@para}{} % Conjugation Paradigm

% Keys
\define@key{ryuhead}{kana}{\def\ryu@head@kana{#1}}
\define@key{ryuhead}{ipa}{\def\ryu@head@ipa{#1}}
\define@key{ryuhead}{pos} {\def\ryu@head@pos{#1}}
\define@key{ryuhead}{cat}{\def\ryu@head@cat{#1}}
\define@key{ryuhead}{acc}{\def\ryu@head@acc{#1}}
\define@key{ryuhead}{rui} {\def\ryu@head@rui{#1}}
\define@key{ryuhead}{para} {\def\ryu@head@para{#1}}

% Keyed Entry head command
\newcommand{\headPrint}[1][]{%
  % reset every time
  \let\ryu@head@kana\@empty%
  \let\ryu@head@ipa\@empty%
  \let\ryu@head@pos\@empty%
  \let\ryu@head@cat\@empty%
  \let\ryu@head@acc\@empty%
  \let\ryu@head@rui\@empty%
  \let\ryu@head@para\@empty%
  % parse keys if any
  \setkeys{ryuhead}{#1}%
  %
  % Kana orthography + header markings
  \ifx\ryu@head@kana\@empty\else%
    \markboth{\ryu@head@kana}{\ryu@head@kana}%
    %\textbf{\ryu@head@kana}%
    {\EmbJpFont\bfseries \ryu@head@kana}%
  \fi%
  % IPA/Romaji orthography
  \ifx\ryu@head@ipa\@empty\else%
    \enspace{\ip{[\ryu@head@ipa]}}%
  \fi%
  % POS block (print as one unit)
  \ifx\ryu@head@pos\@empty\else%
    \kern-0.1em［\ryu@head@pos%
    \ifx\ryu@head@rui\@empty\else\ip{\ryu@head@rui}\fi%
    ］\kern-0.1em%
  \fi%
  % Accent class
  \ifx\ryu@head@acc\@empty\else%
    \kern-0.1em［\ryu@head@acc］\kern-0.1em%
  \fi%
  % Semantic category
  \ifx\ryu@head@cat\@empty\else%
    \kern-0.1em［\ryu@head@cat］\kern-0.1em%
  \fi%
  % Paradigm
  \ifx\ryu@head@para\@empty\else%
    \kern-0.1em［{\EmbJpFont \ryu@head@para}］\kern-0.1em%
  \fi%  
  }%
\makeatother

% Meaning(s) block
\makeatletter
\newcount\ryu@meaningcount
\newcount\ryu@meaningindex

% 2 passes required for the confitioning of numbering 
% i.e. no numbering if only one meaning in meaning block
\newcommand{\MeaningsBlock}[1]{%
  \begingroup%
    % pass 1: count
    \ryu@meaningcount=0\relax%
    \long\def\ryumeaning##1##2##3{%
      \advance\ryu@meaningcount by 1\relax%
    }%
    #1%
    %
    % pass 2: print
    \ryu@meaningindex=0\relax%
    \long\def\ryumeaning##1##2##3{%
      \ifnum\ryu@meaningcount>1\relax%
        \advance\ryu@meaningindex by 1\relax%
        \numbering{\the\ryu@meaningindex}%
      \fi%
      \MeaningPrint{##1}{##2}{##3}%
    }%
    \ignorespaces#1%
  \endgroup%
}
\makeatother

\makeatletter
% Command for printing the meaning
\newcommand{\MeaningPrint}[3]{%
	\ifblank{#1}{}{#1。}% gloss (相当語・訳語)
	\ifblank{#2}{}{#2。}% Explanation ()
	\ifblank{#3}{}{#3}% Example Sentence
}
\makeatother

% Entry Tail
\makeatletter

\newcommand{\ryu@tail@note}{}
\newcommand{\ryu@tail@allo}{}
\newcommand{\ryu@tail@syn}{}

\define@key{ryutail}{note}{\def\ryu@tail@note{#1}}
\define@key{ryutail}{allo}{\def\ryu@tail@allo{#1}}
\define@key{ryutail}{syn} {\def\ryu@tail@syn{#1}}

\newcommand{\tailPrint}[1][]{%
  % reset every time
  \let\ryu@tail@note\@empty%
  \let\ryu@tail@allo\@empty%
  \let\ryu@tail@syn\@empty%
  %
  % parse keys if any
  \setkeys{ryutail}{#1}%
  %
  % output only if nonempty
  \ifx\ryu@tail@note\@empty\else［備］\ryu@tail@note。\fi%
  \ifx\ryu@tail@allo\@empty\else［同］\ryu@tail@allo。\fi%
  \ifx\ryu@tail@syn \@empty\else［類］\ryu@tail@syn 。\fi%
}

\makeatother

% Command for inline sentence example
\newcommand{\inlineExample}[1]{#1。}

% Command for sentence example
% 1st param: Source language orthography
% 2nd param: Target language (Japanese) translation 
\newcommand{\reibun}[2]{\emb{#1}（#2）。}

% Definition of reverse lookup structure
% 1st para: Reading (読み)
% 2nd para: Kanji Orthography (漢字表記)
% 3rd para: Word-form (語形)

% Structure command for reverse lookup
\newcommand{\rentry}[3]{%
  \par\begingroup
  \leavevmode
  \markboth{\unexpanded{#1}}{\unexpanded{#1}}%
  \hangindent=1em
  \noindent
  \ifblank{#2}%
    {#1\hspace{0.4em}}%
    {#1【#2】\hspace{0.4em}}%
  #3%
  \par\endgroup
}
